Alarm Sensor
We created a project with Arduino Uno to display how far an object is from the sensor on an LCD 1602 Module. When an object is close to the sensor, an alarm will sound and a light will turn on.

Inspiration
We were inspired by the Arduino's simplicity and ability to be applied in different situations, especially in ways that can improve our quality of lives.

What it does
We engineered it to display the distance an object, such as a wall, is from the sensor. When the sensor is too close to an object, a continuous sound will play and a light will turn on.

How we built it
First, we put an LCD 1602 Module on the breadboard. Then, we connected jumper wires from the LCD into the respective pins on the Arduino. Next, we used the jumper wire to connect the five-volt slot on the Arduino to the positive of the breadboard. Then, we connected the ground slot on the Arduino to the ground of the breadboard. Then, we connected the anode of the LED to the analog slots on the Arduino and connected the cathode of the LED to a 220 ohm resistor. This was connected to ground of the breadboard. The same process was done for the yellow LED, but to a different analog slot of the Arduino. Lastly, we connected the positive slot of the buzzer to one of the numbered slots of the Arduino (8), and connected the minus slot of the buzzer to ground of the breadboard.

Challenges we ran into
We faced challenges due to the unprecedented complexity of the wiring and the coding. It was our first time with Arduino hardware and software, so we ran into challenges in the beginning, but later in the day, we began to understand how to manipulate code and wires to do what we wanted it to.

Accomplishments that we're proud of
In the beginning, our first milestone was when we got the LCD 1602 Module to display words on its screen. We built off of what we had, and eventually implemented the distance sensor so that the module would display the distance the sensor was from an object in front of it. As we continued to experiment, we were able to add an alarm system and modified it so that if the distance sensed is less than 15 inches, a yellow light will turn on, and if the distance sensed is less than 5 inches, a red light will turn on and a noise will play.

What we learned
We learned that Arduinos are definitely difficult to comprehend! But we overcame that difficulty and worked tirelessly to understand it. We only scratched the surface of Arduinos, so we look forward to what else we can do with them. We learned a lot about the code behind Arduinos and took a lot of time to comprehend what each line meant as well.

What's next for the Alarm Sensor
We will likely continue this project and develop it into something more improved and can be used for various applications.
